\begin{lemma}[2.7.1.1]
\label{IV.2.7.1.1-fr}
Sean $A$ un anillo, $B$ una $A$-álgebra, $A'$ una $A$-álgebra que sea un
$A$-módulo fielmente plano, $B'=B\otimes_A A'$. Para que $B$ sea una
$A$-álgebra de tipo finito \textup{(resp.} de presentación finita\textup{)}, es necesario y
suficiente que $B'$ sea una $A'$-álgebra de tipo finito \textup{(resp.} de presentación
finita\textup{)}.
\end{lemma}

Ya se sabe que la condición es necesaria sin hipótesis sobre $A'$ (1.3.4,
1.3.6, 1.4.3 y 1.4.6). Supongamos que $B'$ es una $A'$-álgebra de tipo finito;
sea $(B_\alpha)_{\alpha\in I}$ la familia filtrante creciente de las
sub-$A$-álgebras de $B$, de modo que se tiene $B=\varinjlim B_\alpha$, y por
tanto también $B'=\varinjlim(B_\alpha\otimes_A A')$, pues el producto tensorial
conmuta con los límites inductivos; si $(x'_i)$ es un sistema finito de
generadores de la $A'$-álgebra $B'$, existe un índice $\alpha$ tal que todos
los $x'_i$ pertenecen a la subálgebra $B_\alpha\otimes_A A'$ de $B'$, de donde
$B'=B_\alpha\otimes_A A'$, y, como $A'$ es fielmente plano,
$B=B_\alpha$ $(\mathbf{0}_{\mathrm I},6.4.1)$.

Supongamos ahora que $B'$ es una $A'$-álgebra de presentación finita; por lo
anterior, ya se sabe que $B$ es una $A$-álgebra de tipo finito, por lo que
existen una $A$-álgebra de polinomios $C=A[T_1,\ldots,T_m]$ y un
$A$-homomorfismo sobreyectivo de álgebras $C\to B$; sea $\mathfrak J$ el
núcleo de este homomorfismo, de modo que se tiene una sucesión exacta
$0\to\mathfrak J\to C\to B\to0$,
y, por tanto, también una sucesión exacta
$0\to\mathfrak J'\to C'\to B'\to0$ (puesto que $A'$ es plano sobre $A$),
poniendo $C'=C\otimes_A A'=A'[T_1,\ldots,T_m]$ y
$\mathfrak J'=\mathfrak J\otimes_A A'$ (identificado con un ideal de $C'$).
Como $B'$ es una $A'$-álgebra de presentación finita, $\mathfrak J'$ es un
$C'$-módulo de tipo finito (1.4.4); pero se tiene
$\mathfrak J'=\mathfrak J\otimes_C C'$ y $C'$ es un $C$-módulo fielmente
plano (2.2.13 y 2.2.3); se sabe entonces que la hipótesis de que $\mathfrak J'$
sea un $C'$-módulo de tipo finito implica que $\mathfrak J$ es un $C$-módulo
de tipo finito (Bourbaki, \emph{Alg. comm.}, cap.~I,
\S~3, núm.~6, prop.~11), por lo que $B$ es una $A$-álgebra de
presentación finita.

(v) se deduce de (iii) y de (2.6.2, (v)) en virtud de (1.5.2).

(vi) se deduce del mismo modo de (iv), (v) y (ii) en virtud de (1.6.1).

(vii) se deduce de (i), (v) y (2.6.4, (ii)) (\textbf{II}, 5.4.1).

(viii) Observemos primero que, como $f'$ es un isomorfismo, es un
homeomorfismo universal, y por tanto lo mismo ocurre con $f$ (2.6.4); se
deduce ya que $f$ es cuasicompacto y separado (2.4.4). Pongamos
$f=(\psi,\theta)$, donde $\psi$ es, pues, un homeomorfismo; se trata de
demostrar que
$\theta:\mathcal O_Y\to f_*(\mathcal O_X)$
es un isomorfismo de $\mathcal O_Y$-módulos. Ahora bien, si se pone
$f'=(\psi',\theta')$, el homomorfismo
$\theta':\mathcal O_{Y'}\to f'_*(\mathcal O_{X'})$ está compuesto por el
homomorfismo canónico
$g^*(f_*(\mathcal O_X))\to f'_*(\mathcal O_{X'})$ y por $g^*(\theta)$
(2.3.2); pero el primero de estos dos homomorfismos es biyectivo en virtud de
la hipótesis sobre $g$ (2.3.1), por lo que, si $\theta'$ es biyectivo,
\oldpage[IV-2]{31}
lo mismo ocurre con $g^*(\theta)$, y, como $g$ es fielmente plano, $\theta$ es
biyectivo (2.2.7), lo que demuestra (viii).

(ix) La proposición se deduce de (viii), de (\textbf I, 5.3.4) y de (\textbf I, 5.3.8),
que reduce los monomorfismos a los isomorfismos.

(x) Si $f'$ es una inmersión abierta, $f'(X')$ es abierto en $Y'$, y
se tiene $f'(X')=g^{-1}(f(X))$ (\textbf I, 3.4.8); de (2.3.12) se deduce que $f(X)$
es abierto. Se puede entonces reemplazar $Y$ (resp. $Y'$) por el subpreesquema
inducido sobre el abierto $f(X)$ (resp. $f'(X')$), teniendo en cuenta (1.1.2)
y (2.2.13); entonces $f'$ se convierte en un isomorfismo, por lo que lo mismo
ocurre con $f$ por (viii), y esto establece (x).

(xi) Si $f'$ es una inmersión cuasicompacta, $f'$ es un morfismo
cuasicompacto y cuasiseparado (1.2.2), por lo que lo mismo ocurre con $f$ por
(ii) y (2.6.2, (v)). Sea $Z$ el subpreesquema de $Y$, imagen cerrada de $X$
por $f$ (1.7.8), y pongamos $f=j\circ q$, donde $j:Z\to Y$ es la inyección
canónica; se tiene entonces $f'=j'\circ q'$, con $j'=j_{(Y')}$,
$q'=q_{(Y')}$, y se sabe que $j'$ se identifica con la inyección canónica
$Z'\to Y'$ del subpreesquema $Z'$ de $Y'$, imagen cerrada de $X'$ por $f'$
(2.3.2). La hipótesis sobre $f'$ significa entonces que $q'$ es una inmersión
abierta (\textbf I, 9.5.10), por lo que lo mismo ocurre con $q$ por (x), y esto muestra
que $f$ es una inmersión.

(xii) Decir que $f$ (resp. $f'$) es una inmersión cerrada significa que
$f$ (resp. $f'$) es una inmersión cuasicompacta y un morfismo cerrado; se ve,
por tanto, que (xii) se deduce de (xi) y de (2.6.2, (ii)).

(xiii) y (xiv) Supongamos que $f'$ es afín (resp. cuasiafín);
observemos entonces que $f'$ es cuasicompacto y cuasiseparado (\textbf{II}, 5.1.1), por
lo que lo mismo ocurre con $f$ por (ii) y (2.6.2, (v)). Pongamos
$\mathcal A=f_*(\mathcal O_X)$, $\mathcal A'=f'_*(\mathcal O_{X'})$; en
virtud de (2.3.1), el homomorfismo canónico de $\mathcal O_{Y'}$-álgebras
$g^*(\mathcal A)\to\mathcal A'$ es biyectivo; por consiguiente, si
$h:Z=\operatorname{Spec}(\mathcal A)\to Y$ es el morfismo estructural, el
morfismo estructural $h':Z'=\operatorname{Spec}(\mathcal A')\to Y'$ se
identifica con $h_{(Y')}$ (\textbf{II}, 1.5.2). Sea entonces $u:X\to Z$ (resp.
$u':X'\to Z'$) el $Y$-morfismo (resp. $Y'$-morfismo) canónico correspondiente
al homomorfismo identidad de $\mathcal A$ (resp. $\mathcal A'$) (\textbf{II}, 1.2.7);
como se tiene el diagrama conmutativo
\[
\xymatrix{
X \ar[d]_{u} & X' \ar[l] \ar[d]^{u'}\\
Z \ar[d]_{h} & Z' \ar[l]_{g'} \ar[d]^{h'}\\
Y & Y' \ar[l]^{g}
}
\]
y $h'\circ u'=f'$, de (\textbf{II}, 1.2.7) se deduce que $u'=u_{(Z')}$. Además, $g'$
es fielmente plano y cuasicompacto (1.1.2 y 2.2.13). Dicho esto, la hipótesis
sobre $f'$ significa que $u'$ es un isomorfismo (resp. una inmersión abierta)
(\textbf{II}, 5.1.6); de (viii) (resp. (x)) se deduce entonces que $u$ es un
isomorfismo (resp. una inmersión abierta), de donde se sigue (xiii) (resp.
(xiv)).

\oldpage[IV-2]{32}
(xv) Si $f'$ es finito, es afín, por lo que también $f$ lo es según
(xiii); además, con las notaciones de la demostración de (xiii), $\mathcal A'$
es un $\mathcal O_{Y'}$-módulo de tipo finito, y $\mathcal A'$ es isomorfo a
$g^*(\mathcal A)$; de (2.5.2) se deduce que $\mathcal A$ es un
$\mathcal O_Y$-módulo de tipo finito, por lo que $f$ es un morfismo finito.

(xvi) Decir que $f$ es cuasifinito significa que $f$ es un morfismo de
tipo finito y que, para todo $y\in Y$, $f^{-1}(y)$ es finito (\textbf{II}, 6.2.2 y \textbf{I},
6.4.4); la conclusión se deduce, por tanto, de (v) y de (xv).

(xvii) Se ve, como en (xv), que $f$ es afín. Es posible limitarse al caso
en que $Y=\operatorname{Spec}(A)$, $Y'=\operatorname{Spec}(A')$, y entonces
$X=\operatorname{Spec}(B)$, $X'=\operatorname{Spec}(B')$, donde
$B'=B\otimes_A A'$; $B$ es igual al límite inductivo de sus
sub-$A$-álgebras de tipo finito $B_\alpha$, por lo que se tiene
$B'=\varinjlim B'_\alpha$, donde $B'_\alpha=B_\alpha\otimes_A A'$, y
$B'_\alpha$ es una $A'$-álgebra de tipo finito. Pero, por hipótesis, $B'$ es
entero sobre $A'$, por lo que $B'_\alpha$ es un $A'$-\emph{módulo} de tipo finito, y
$B_\alpha$ es entonces un $A$-módulo de tipo finito (2.5.2). Q.E.D.
